\chapter{Maps, Folds, and the Combinator-Primary Model}
\label{ch:combinators}

\chapterorient{The previous chapter built one combinator, \texttt{iterate\_while},
and showed that every iteration in the book is a setting of it. This chapter widens
the lens. \texttt{iterate\_while} is the iteration primitive; it has \emph{relatives}
that organize the rest of the simulator---\code{map} for independent work,
\code{fold} for chained work, and a small algebra of data-combining verbs the steps
call. We give the four solve corners of \cref{ch:situating} their combinator names,
we meet the tensor-sum combinator \code{linear\_combination} and watch the familiar
BLAS routines fall out of it as specializations, and we state the chapter's
thesis---the \emph{combinator-primary model}: a few higher-order combinators are
the load-bearing vocabulary, and the many named operators are \emph{specialization
notes} hung beneath them. This is why a book about a hundred-thousand-line simulator
can be short.}

\section{From one combinator to a family of relatives}

\Cref{ch:fold} earned a strong claim: there is \emph{one} iteration combinator,
\code{iterate\_while}, and independent steps, sequential carries, and convergence
loops are all that combinator at different settings. We will not retreat from it.
But ``one combinator expresses every iteration'' is a statement about
\emph{control flow}---about how a loop threads its carry and tests its predicate. A
simulator is more than its loops. It also \emph{combines data}: it sums scaled
tensors, it forms inner products, it assembles operators. And when it does run a
family of solves, the \emph{shape} of that family---independent versus
chained---deserves a name of its own, because the two shapes have radically
different consequences for parallelism.

So this chapter steps up one level. Where \cref{ch:fold} looked \emph{inside} a
single loop, this chapter looks at the small set of combinators that organize the
\emph{whole} simulator: the ones that drive families of solves, and the ones that
combine tensors. \Cref{fig:relatives} lays out the family tree. At the root is the
distinction the whole chapter turns on---\code{map} versus \code{fold}, independent
versus carried---and beneath it hang the named combinators a driver actually calls.

\begin{figure}[t]
\centering
\begin{tikzpicture}[font=\footnotesize, node distance=4mm]
  % root
  \node[stage, minimum width=58mm, align=center] (root) at (0,0)
    {the combinator vocabulary};
  % two branches: iteration-structural vs data-algebra
  \node[stage, minimum width=44mm, align=center, fill=simBgGreen, draw=simGreen]
    (iter) at (-3.6,-1.7) {iteration-structural\\(thread a carry / family)};
  \node[stage, minimum width=44mm, align=center, fill=simBgGold, draw=simGold]
    (data) at (3.6,-1.7) {data-algebra\\(combine tensors)};
  \draw[depedge] (root) -- (iter);
  \draw[depedge] (root) -- (data);
  % iteration-structural leaves
  \node[opbox, minimum width=40mm, align=left, below=10mm of iter] (il) {%
    \code{solve\_family} \;\itshape map (independent)\\
    \code{frequency\_sweep} \;\itshape op-varying map\\
    \code{fold\_solve} \;\itshape fold (chained)\\
    \code{iterate\_while} \;\itshape the loop primitive};
  \draw[depedge] (iter) -- (il);
  % data-algebra leaves
  \node[opbox, minimum width=40mm, align=left, below=10mm of data] (dl) {%
    \code{linear\_combination} \;\itshape tensor sum\\
    \quad \code{scal}/\code{axpy}/\code{axpby}/\dots \;\itshape leaves\\
    \code{inner\_product} \;\itshape scalar fold\\
    \quad \code{dot}/\code{nrm2} \;\itshape leaves};
  \draw[depedge] (data) -- (dl);
\end{tikzpicture}
\caption[The combinator vocabulary]{The combinator vocabulary of the simulator,
in two branches. The \emph{iteration-structural} combinators (left) thread a carry
or drive a family of solves; they are relatives of the \code{iterate\_while} of
\cref{ch:fold}, distinguished by whether the family is an independent \emph{map}
or a chained \emph{fold}. The \emph{data-algebra} combinators (right) combine
tensors into tensors or scalars; each has a cluster of named BLAS-style leaves
hanging beneath it. The whole of this short chapter is a tour of these two boxes
and the one thesis they illustrate: the combinator is the entry, the leaves are
notes.}
\label{fig:relatives}
\end{figure}

\section{The two shapes of a family: map and fold}
\label{sec:mapfold}

Before naming the simulator's combinators, fix the one distinction everything
rests on. We met it in \cref{ch:fold} as the contrast between \code{map} and
\code{foldl}; here we make it the organizing axis.

A \code{map} applies the \emph{same} operation to each element of a collection
\emph{independently}:
\begin{lstlisting}[style=l4style]
map f [x1, x2, x3] = [f x1, f x2, f x3]
\end{lstlisting}
The three results do not depend on one another. \code{f x2} can be computed before,
after, or at the same time as \code{f x1}; you could run each on a separate
machine and collect the answers in any order. A \code{map} has \emph{no carry}.

A \code{foldl} threads an accumulator---a \emph{carry}---through the collection,
each step consuming what the last produced:
\begin{lstlisting}[style=l4style]
foldl step s0 [x1, x2, x3] = step (step (step s0 x1) x2) x3
\end{lstlisting}
Here the order is forced. The accumulator entering the third step is exactly what
the second step produced; you cannot start step three until step two has finished.
A \code{fold} is a \emph{chain}.

This is the single most consequential distinction in the simulator's architecture,
and \cref{fig:mapfold} draws it bare. The map fans out; the fold lines up. The map
parallelizes; the fold serializes. Everything else in this section is a
specialization of one or the other.

\begin{figure}[t]
\centering
\begin{tikzpicture}[font=\footnotesize]
  % ===== LEFT: map = independent fan-out =====
  \begin{scope}
    \node[font=\small\bfseries, text=simBlue] at (1.4,2.7) {a \code{map}: independent fan-out};
    \node[data, minimum width=10mm] (src) at (1.4,1.9) {family};
    \foreach \i/\x in {1/0,2/1.4,3/2.8}{
      \node[opbox, minimum width=9mm] (m\i) at (\x,0.8) {\code{f}};
      \node[product, minimum width=9mm, minimum height=5mm] (r\i) at (\x,-0.4) {$y_{\i}$};
      \draw[flow] (src.south) -- (m\i.north);
      \draw[flow] (m\i.south) -- (r\i.north);
    }
    \node[font=\scriptsize, text=simGreen, align=center] at (1.4,-1.3)
      {no arrows \emph{between} the \code{f}'s:\\reorderable, parallelizable};
  \end{scope}
  \draw[simGray!50, thick] (4.2,-1.6) -- (4.2,3.0);
  % ===== RIGHT: fold = sequential chain =====
  \begin{scope}[xshift=5.0cm]
    \node[font=\small\bfseries, text=simGreen] at (2.4,2.7) {a \code{fold}: sequential chain};
    \node[data, minimum width=9mm] (s0) at (-0.3,0.8) {$s_0$};
    \foreach \i/\x in {1/1.1,2/2.5,3/3.9}{
      \node[opbox, minimum width=9mm] (g\i) at (\x,0.8) {\code{step}};
    }
    \node[data, minimum width=9mm] (s3) at (5.3,0.8) {$s_3$};
    \draw[flow] (s0.east) -- (g1.west);
    \draw[flow] (g1.east) -- node[above,font=\scriptsize]{$s_1$} (g2.west);
    \draw[flow] (g2.east) -- node[above,font=\scriptsize]{$s_2$} (g3.west);
    \draw[flow] (g3.east) -- (s3.west);
    \foreach \i/\x in {1/1.1,2/2.5,3/3.9}{
      \node[font=\scriptsize, text=simGold!80!black] at (\x,1.65) {$x_{\i}$};
      \draw[refedge, simGold!80!black] (\x,1.5) -- (g\i.north);
    }
    \node[font=\scriptsize, text=simRust, align=center] at (2.5,-0.6)
      {each \code{step} needs the previous $s$:\\strictly ordered, not parallelizable};
  \end{scope}
\end{tikzpicture}
\caption[Map versus fold: the fundamental split]{The fundamental split. A
\emph{map} (left) applies \code{f} to each family element independently---the
three applications share no data, so they fan out, reorder freely, and run in
parallel. A \emph{fold} (right) threads a carry $s_0 \to s_1 \to s_2 \to s_3$
through the chain---each \code{step} consumes the previous carry, so the chain is
strictly ordered and cannot be parallelized. The presence or absence of the
horizontal carry arrows is the entire difference, and it decides whether the work
is embarrassingly parallel or intrinsically sequential.}
\label{fig:mapfold}
\end{figure}

\begin{keyidea}
A \emph{map} runs a family of \emph{independent} computations: no carry, so the
elements commute, reorder, and parallelize. A \emph{fold} runs a \emph{chain}:
each step consumes the previous step's output, so the steps are strictly ordered
and intrinsically sequential. This one axis---is there a carry between the
elements?---decides the parallelism of every family of solves the simulator runs.
A map is in fact the degenerate fold whose step \emph{ignores} the carry; the two
are one family, split by whether the carry is used.
\end{keyidea}

The claim that a map is a degenerate fold is worth making exact, because it is the
formal hinge on which the whole ``one combinator'' story of \cref{ch:fold} turns.
A \code{foldl} that drops its accumulator on every step does no threading at all:
the result of each step depends only on the element, not on the carry, so the steps
are independent and we have a map in disguise.

\begin{proposition}[Map is the carry-dropping fold]
\label{prop:mapfold}
For any element function \code{f}, the map \lstinline[style=l4style]|map f xs| equals
the trajectory of a left fold whose step ignores its accumulator and emits
\code{f} of the element:
\begin{lstlisting}[style=l4style]
map f xs == (foldl_collecting (\_ x -> f x) seed xs).collected
\end{lstlisting}
where \code{foldl\_collecting} threads a carry but \emph{also} appends each step's
output to a collected list. Because the step's first argument (the carry) never
appears on its right-hand side, the steps share no data, and the collected list is
order-independent.
\end{proposition}

The proposition is not deep---it is almost a restatement of the definitions---but
its \emph{consequence} is the architecture of the simulator. It says the map and
the fold are not two unrelated control structures to be learned separately; they
are one structure, the value-threading fold, observed at two settings of a single
switch: \emph{does the step read its carry?} When the answer is no, the work fans
out and parallelizes; when the answer is yes, the work lines up and serializes. The
entire taxonomy of the next section is this one switch, crossed with a second
(does the operator change?).

It pays to draw out the three practical consequences of the split, because each
recurs in every later Part.

\begin{description}[style=nextline, leftmargin=0pt, font=\normalfont\itshape]
  \item[Parallelism follows the carry.] A map's elements may be evaluated in any
  order, on any number of machines, and reassembled; a fold's steps must run in
  sequence on one timeline. When \cref{part:modalities} reports that electrostatics
  ``scales trivially across conductors'' but the transient analysis ``cannot
  overlap its timesteps,'' it is reporting exactly this---the first is a map, the
  second a fold.
  \item[Reordering is a correctness question, not a performance one.] For a map,
  reordering the family permutes the output identically and changes nothing else;
  for a fold, reordering the steps changes the \emph{answer}, because step $k$'s
  input is step $k{-}1$'s output. ``Can I reorder this?'' is therefore answered by
  ``is there a carry?''---a structural fact visible in the combinator's type, not a
  property to be discovered by testing.
  \item[Memory follows the carry too.] A map need keep only the inputs and the
  collected outputs; a fold must additionally keep the live carry as it threads.
  When the carry is a whole field state (the transient analysis) this is the
  dominant cost, and \cref{ch:transient} returns to it. A map has no such live
  intermediate.
\end{description}

\section{The four solve corners, as combinators}
\label{sec:corners}

We can now return to the four solve corners of \cref{ch:situating} and give each
its combinator name. There we classified a driver's solve stage by two questions:
is the operator \emph{fixed} or \emph{varying} as we proceed, and is the
computation a \emph{map} or a \emph{fold}? Crossing those questions gave a
$2\times2$ grid with three occupied cells and a fourth corner sitting outside it.
\Cref{fig:cornercomb} redraws that grid, now labelling each cell with the
combinator the framework uses for it.

\begin{figure}[t]
\centering
\begin{tikzpicture}[font=\footnotesize]
  \node[cornertag, anchor=south] at (1.2,4.5) {operator \emph{fixed}};
  \node[cornertag, anchor=south] at (5.6,4.5) {operator \emph{varying}};
  \node[cornertag, rotate=90, anchor=south] at (-1.5,3.2) {a \textbf{map}};
  \node[cornertag, rotate=90, anchor=south] at (-1.5,0.9) {a \textbf{fold}};
  % top-left: solve_family
  \node[stage, minimum width=34mm, minimum height=20mm, anchor=south west,
    fill=simBgBlue, align=center] (c1) at (0,2.4) {};
  \node[anchor=north, font=\scriptsize\bfseries] at (1.7,4.1) {fixed-operator map};
  \node[anchor=center, font=\small\ttfamily, text=simBlue] at (1.7,3.3) {solve\_family};
  \node[anchor=north, font=\scriptsize\itshape] at (1.7,2.95) {electrostatic,};
  \node[anchor=north, font=\scriptsize\itshape] at (1.7,2.65) {magnetostatic};
  % top-right: frequency_sweep
  \node[stage, minimum width=34mm, minimum height=20mm, anchor=south west,
    fill=simBgBlue, align=center] (c2) at (4.4,2.4) {};
  \node[anchor=north, font=\scriptsize\bfseries] at (6.1,4.1) {operator-varying map};
  \node[anchor=center, font=\small\ttfamily, text=simBlue] at (6.1,3.3) {frequency\_sweep};
  \node[anchor=north, font=\scriptsize\itshape] at (6.1,2.85) {driven};
  % bottom-left: fold_solve
  \node[stage, minimum width=34mm, minimum height=20mm, anchor=south west,
    fill=simBgGreen, align=center] (c3) at (0,0) {};
  \node[anchor=north, font=\scriptsize\bfseries] at (1.7,1.7) {state-threaded fold};
  \node[anchor=center, font=\small\ttfamily, text=simGreen] at (1.7,0.95) {fold\_solve};
  \node[anchor=north, font=\scriptsize\itshape] at (1.7,0.55) {transient, AMR loop};
  % bottom-right: eigsolve (outside)
  \node[extern, minimum width=34mm, minimum height=20mm, anchor=south west,
    align=center] (c4) at (4.4,0) {};
  \node[anchor=north, font=\scriptsize\bfseries] at (6.1,1.7) {opaque black-box};
  \node[anchor=center, font=\small\ttfamily, text=simViolet] at (6.1,0.95) {eigsolve};
  \node[anchor=north, font=\scriptsize\itshape] at (6.1,0.55) {eigenmode, boundary-mode};
  \node[font=\scriptsize\itshape, text=simViolet] at (6.1,-0.35) {(no loop we write)};
\end{tikzpicture}
\caption[The four solve corners as combinators]{The four solve corners of
\cref{ch:situating}, now named as combinators (compare \cref{fig:corners}). A
fixed-operator map is \code{solve\_family}: assemble once, map independent solves.
An operator-varying map is \code{frequency\_sweep}: rebuild the operator inside the
map. A state-threaded fold is \code{fold\_solve}: chain the solves, each from the
last. The fourth corner is \code{eigsolve}, the opaque library call that sits
outside the map/fold grid because we write no loop at all. Three of the four are
relatives of \code{iterate\_while}; the fourth is the deliberate non-loop.}
\label{fig:cornercomb}
\end{figure}

\subsection{\texttt{solve\_family}: the fixed-operator map}

The first corner is a \code{map}. Assemble the operator $\Kop$ \emph{once}, then
solve $\Kop\,\vx_i = \vb_i$ for a family of right-hand sides $\vb_1,\dots,\vb_n$.
The operator never changes; only the right-hand side does. Because the solves
share one operator and depend on nothing but their own right-hand side, they are
independent: a \code{map}. The combinator is \code{solve\_family}, and its
definition is a one-liner---a \code{map} of the single-solve cap over the family:
\begin{lstlisting}[style=l4style]
solve_family :: OpParams -> [Inputs] -> [SimState]
solve_family op rhss = map (\inp -> ksp_solve op inp) rhss
\end{lstlisting}
The load-bearing fact is the \emph{hoist}: \code{op} is captured once, outside the
\code{map}, and threaded unchanged into every per-element solve. The solver built
from \code{op}---a factorization, a preconditioner---is invariant across the map,
so it too is built once and reused. This is exactly the Palace pattern where the
operator setup call sits \emph{outside} the sweep loop: electrostatics and
magnetostatics both assemble $\Kop$ once and then loop over conductors, reusing the
one operator. The exemplars are electrostatics (\cref{ch:electrostatics}) and
magnetostatics, which are this corner with different operators and different
families.

Because the map carries nothing between elements, \code{solve\_family} obeys the
\emph{concatenation law}: \lstinline[style=l4style]|solve_family op (a ++ b) = solve_family op a ++ solve_family op b|.
Splitting the family, chunking it, reordering it, or solving its pieces on separate
machines all give the same answer---the algebraic statement of embarrassing
parallelism. There is a companion \emph{order law}:
\lstinline[style=l4style]|solve_family op (permute rhss) = permute (solve_family op rhss)|,
which says the underlying solves commute even though the collection preserves
position. Together the two laws are the formal license for every parallel
realization of the fixed-operator corner: the family may be partitioned arbitrarily,
solved out of order, and stitched back together, and the answer is unchanged. Note
what the laws rest on---a \emph{shared, unchanging} \code{op}. Strip that away and
the next combinator forfeits both laws.

\subsection{\texttt{frequency\_sweep}: the operator-varying map}

Keep the map, but change one thing: the operator is no longer fixed. The driven
analysis (\cref{ch:driven}) computes scattering parameters by solving the
time-harmonic system at each frequency $\omega$ in a swept band, and the system
operator \emph{depends on frequency}:
\[
  \mat{A}(\omega) \;=\; \Kop + i\omega\,\Cop - \omega^2\Mop.
\]
So before each solve, inside the sweep, the operator must be rebuilt at the current
$\omega$. The solves are still independent---the answer at $\omega_3$ does not need
the answer at $\omega_2$---so this is still a \code{map}, not a fold. What
distinguishes it from \code{solve\_family} is precisely that the operator is rebuilt
\emph{inside} the loop rather than hoisted outside it. The combinator is
\code{frequency\_sweep}:
\begin{lstlisting}[style=l4style]
frequency_sweep :: OperatorFamily -> [Omega] -> [SimState]
frequency_sweep fam omegas =
  map (\w -> let aw = assemble_frequency_operator fam w   -- REBUILD inside the map
             in  ksp_solve aw (rhs_at w))
      omegas
\end{lstlisting}
The basis $\{\Kop,\Cop,\Mop\}$ is captured once; the per-member operator
$\mat{A}(\omega)$ is rebuilt per member. The contrast with \code{solve\_family} is
a single axis---is the operator captured once and hoisted, or rebuilt per
element?---and it is the entire structural difference between electrostatics and
the driven analysis.

\begin{pitfall}
\code{frequency\_sweep} and \code{solve\_family} differ by the \emph{position of
one setup call} relative to the map: outside the loop (hoist) versus inside it
(rebuild). This is not an optimization knob. An implementation that ``speeds up''
\code{frequency\_sweep} by hoisting the operator construction out of the map does
not run faster---it computes the \emph{wrong answer}, because it solves every
frequency against the operator of the first one. The hoist is correct exactly when
the operator is fixed and a bug exactly when it varies. The combinator names which
case you are in.
\end{pitfall}

\subsection{\texttt{fold\_solve}: the state-threaded fold}

Now drop the independence. In \code{solve\_family} and \code{frequency\_sweep} the
solves were a map: independent, reorderable, parallelizable. In the third corner
they are a \code{fold}: each step begins from the state the previous step produced.
The transient analysis (\cref{ch:transient}) marches forward in time---the field
state at $t_{k+1}$ is computed from the state at $t_k$ by one time step---and the
state is \emph{threaded} through the chain. The combinator is \code{fold\_solve},
a \code{foldl} of an opaque per-step operator over a schedule:
\begin{lstlisting}[style=l4style]
fold_solve :: OpParams -> TimeState -> [Time] -> TimeState
fold_solve op s0 schedule = foldl (\s t -> time_step_op op s t) s0 schedule
\end{lstlisting}
The signature alone tells the story. The seed \code{s0} and the result are both a
\code{TimeState}; the step \lstinline[style=l4style]|time_step_op op s t| reads
\code{s}, the prior carry. That read is what makes this a fold and not a map: it
forbids reordering at the type level. The adaptive-mesh-refinement outer loop of
\cref{ch:situating} is also a \code{fold\_solve} of this shape---each refinement
pass begins from the mesh the last pass produced---which is why the same combinator
serves both transient time-stepping and the AMR loop.

Two things are worth pausing on. First, the per-step operator
\lstinline[style=l4style]|time_step_op| is, for the transient analysis,
\emph{opaque}: it bottoms out in a library time-integrator (an implicit Runge--Kutta
or generalized-$\alpha$ step) that the combinator only quantifies over, rather than
rendering. \code{fold\_solve} names the fold structure and the carry-threading; what
one step \emph{does} is a separate concern, deferred to \cref{ch:transient}. The
combinator is honest about the boundary the way \code{eigsolve} is, but it keeps the
loop in our hands---we own the fold, the library owns the step.

Second, the laws \emph{invert} relative to the map. \code{fold\_solve} obeys
neither the concatenation law nor the order law: splitting the schedule and folding
the pieces independently is wrong (the second piece must start from where the first
left off, not from \code{s0}), and permuting the schedule changes the answer
outright. What it obeys instead is a \emph{decomposition} law---folding a
concatenated schedule equals folding the first part and then folding the second
\emph{from the carry the first produced}:
\begin{lstlisting}[style=l4style]
fold_solve op s0 (a ++ b) = fold_solve op (fold_solve op s0 a) b
\end{lstlisting}
This is the algebraic signature of sequentiality: the only way to combine two folds
is to thread the carry from one into the other. Compare \code{solve\_family}, whose
concatenation law combines two maps by simply \emph{appending} their outputs with no
threading at all. The two laws---append versus thread-through---are the map/fold
split written algebraically, and they are why a map distributes over a cluster while
a fold marches on a single timeline.

\subsection{\texttt{eigsolve}: the corner that is not a loop}

The fourth corner is different in kind. In the first three \emph{we} write the
loop---the map or the fold---and can see every solve it performs. In the fourth we
write no loop at all. The eigenmode analysis (\cref{ch:eigenvalues}) assembles its
operator pencil once and hands the whole problem to a specialized eigenvalue
library (SLEPc, in Palace's case) with a single call. The iteration that produces
the answer runs \emph{inside} the library, where our framework cannot see it. The
combinator \code{eigsolve} is therefore deliberately \emph{opaque}: it names the
library boundary and the contract across it and stops. It is not a relative of
\code{iterate\_while}; it is the marked absence of a loop we own. (\Cref{ch:eigenvalues}
separately reconstructs what the library does internally, in our own vocabulary,
turning the black box into a glass box for the curious reader---but the
\emph{combinator} that drives the eigenmode analysis is just the single opaque
call.)

\begin{keyidea}
The four solve corners \emph{are} four combinators. \code{solve\_family} is the
fixed-operator map (assemble once, map independent solves). \code{frequency\_sweep}
is the operator-varying map (rebuild the operator inside the map).
\code{fold\_solve} is the state-threaded fold (chain the solves, each from the
last). \code{eigsolve} is the opaque black-box (no loop we write). Three are
relatives of \code{iterate\_while}, split by the map/fold axis and the
fixed/varying axis; the fourth is the deliberate non-loop. Every driver in
\cref{part:modalities} is one of these four over the \texttt{assemble} $\to$
\texttt{solve} $\to$ \texttt{reduce} skeleton.
\end{keyidea}

\section{The data-algebra combinators}
\label{sec:dataalgebra}

The solve-coordination combinators of \cref{sec:corners} say \emph{how} a family of
solves is structured. Inside each solve, the step body \emph{combines tensors}:
it scales a search direction and adds it to the iterate, it forms the inner product
that computes a step length, it measures a residual norm. These combinations are
the \emph{data-algebra} combinators, the right-hand branch of \cref{fig:relatives}.
There are only two that matter, and each illustrates the chapter's thesis in
miniature: one combinator, several named specializations.

\subsection{\texttt{linear\_combination}: the tensor-sum combinator}

The first is \code{linear\_combination}, the scalar-weighted tensor sum. Given a
list of (coefficient, tensor) pairs, it forms $\sum_i a_i\,t_i$:
\begin{lstlisting}[style=l4style]
linear_combination :: [(Scalar, Tensor)] -> Tensor
linear_combination pairs = foldl (\acc (a, t) -> acc + scal a t) zeros pairs
\end{lstlisting}
It is a \code{foldl} over the term list, accumulating $a_i\,t_i$ into a running
tensor that starts at zero. Every term must be \emph{congruent}---the same shape---
because the sum is element-local; the result has that shape too. This is the
single most-used data combinator in the book: a Krylov step's update
$\vx \gets \vx + \alpha\,\vp$ is a \code{linear\_combination}; GMRES's basis
correction $\sum_j y_j\,\vv_j$ is a \code{linear\_combination}; the driven
operator $\Kop + i\omega\Cop - \omega^2\Mop$ is a \code{linear\_combination} of
\emph{operators} rather than vectors.

It satisfies the laws you would hope for, and the one that earns its keep is the
\emph{concatenation-homomorphism}:
\[
  \texttt{linear\_combination}\,(a \mathbin{+\!\!+} b)
    \;=\; \texttt{linear\_combination}\,a \;+\; \texttt{linear\_combination}\,b.
\]
Splitting the term list and summing the pieces gives the same tensor as summing the
whole. In the language of \cref{ch:language} this makes \code{linear\_combination} a
\emph{monoid homomorphism}: it carries the monoid of term lists (under
concatenation, with the empty list as unit) to the monoid of tensors (under
addition, with the zero tensor as unit). That single structural fact organizes the
rest of the verb's algebra. The empty list maps to the unit:
\lstinline[style=l4style]|linear_combination [] = zeros|. A zero coefficient drops
its term: \lstinline[style=l4style]|linear_combination ((0,t):rest) = linear_combination rest|
(the algebraic content of the implementation's ``skip the term if its coefficient is
zero'' branch). A scalar can move from the coefficient onto the tensor:
$\code{linear\_combination}\,((\kappa a, t):\text{rest}) = \code{linear\_combination}\,((a, \kappa t):\text{rest})$.
And over exact arithmetic the sum is permutation-invariant---reordering the term
list leaves the result unchanged---because tensor addition is commutative and
associative. (Under floating-point arithmetic that last law holds only
approximately, since summation order perturbs the rounding; the framework treats
the left-to-right fold order as canonical and records the reordering subtlety as a
load-bearing numerical note rather than an exact law, in the spirit of
\cref{ch:bigidea}'s distinction between transparent and load-bearing tricks.)

The concatenation-homomorphism is also the law that makes the four named BLAS
routines \emph{one} combinator, as we see next: it lets a longer fixed-arity sum be
built from shorter ones, which is exactly the subsumption chain of the arity leaves.

\subsection{The BLAS verbs are arity leaves}

Numerical libraries ship a family of fused routines for small fixed-length sums,
the BLAS-1 ``\code{axpy}'' family. Each is \code{linear\_combination} at one fixed
term-list length:
\begin{lstlisting}[style=l4style]
scal(a, x)                 = linear_combination [(a, x)]                  -- arity 1
axpy(a, x, y)              = linear_combination [(a, x), (1, y)]          -- arity 2, 2nd coeff = 1
axpby(a, x, b, y)          = linear_combination [(a, x), (b, y)]          -- arity 2, general
axpbypcz(a, x, b, y, c, z) = linear_combination [(a, x), (b, y), (c, z)]  -- arity 3
\end{lstlisting}
Read top to bottom: \code{scal} scales one tensor (a one-term list); \code{axpy}
computes $a\vx + \vy$ (a two-term list whose second coefficient is fixed to $1$);
\code{axpby} computes $a\vx + b\vy$ (the general two-term list); \code{axpbypcz}
computes $a\vx + b\vy + c\vz$ (the three-term list). They differ \emph{only} in the
length of the term list. The concatenation law makes the subsumption exact:
\code{axpbypcz}'s three-term list is the concatenation of \code{axpby}'s two-term
list with \code{scal}'s one-term list, so $\code{scal} \prec \code{axpy} \prec
\code{axpby} \prec \code{axpbypcz}$ is the bounded-arity shadow of the one
combinator.

These fused routines exist because a hand-written kernel for ``$a\vx + b\vy$''
touches memory once instead of twice and so runs faster on a CPU. That is a real
performance win, but it is \emph{not} an abstraction: \code{axpby} adds no meaning
that \code{linear\_combination} lacks. So the framework makes a deliberate choice.
The accelerated kernels are ``stopped low''---kept as performance notes at the
implementation level---while the \emph{combinator rises} to the top of the
vocabulary in their place. \Cref{fig:lincomb} draws this: \code{linear\_combination}
is the trunk; the four BLAS routines hang off it as arity leaves, labels for fixed
list-lengths, not combinators in their own right.

The same combinator generalizes past vectors. Its terms need not be column tensors;
they may be \emph{operators}. The driven analysis's frequency-dependent operator
$\mat{A}(\omega) = \Kop + i\omega\,\Cop - \omega^2\Mop$ is a
\code{linear\_combination} of the three assembled operators $\Kop, \Cop, \Mop$ with
the $\omega$-dependent coefficients $1, i\omega, -\omega^2$. This is the same fold,
the same concatenation law, the same arity story---only the term type has changed
from tensor to operator. When \cref{ch:driven} writes ``rebuild $\mat{A}(\omega)$''
inside the \code{frequency\_sweep} map, the rebuild it means is one
\code{linear\_combination} over operator terms. The two combinator branches of
\cref{fig:relatives} meet here: the operator-varying map calls the tensor-sum
combinator to do its per-step rebuild. Recognizing the operator assembly as ``the
same verb, operator-valued terms'' rather than a new primitive is the
combinator-primary model paying off inside the framework itself.

\begin{figure}[t]
\centering
\begin{tikzpicture}[font=\footnotesize, node distance=5mm]
  % trunk
  \node[stage, minimum width=64mm, align=center] (trunk) at (0,0)
    {\code{linear\_combination} \;\itshape\quad $\sum_i a_i\,t_i$ \;\;(term list, any length)};
  % four leaves
  \def\ly{-1.9}
  \node[opbox, minimum width=24mm, align=center] (l1) at (-4.8,\ly)
    {\code{scal}\\\scriptsize arity 1\\\scriptsize $a x$};
  \node[opbox, minimum width=24mm, align=center] (l2) at (-1.6,\ly)
    {\code{axpy}\\\scriptsize arity 2\\\scriptsize $a x + y$};
  \node[opbox, minimum width=24mm, align=center] (l3) at (1.6,\ly)
    {\code{axpby}\\\scriptsize arity 2\\\scriptsize $a x + b y$};
  \node[opbox, minimum width=24mm, align=center] (l4) at (4.8,\ly)
    {\code{axpbypcz}\\\scriptsize arity 3\\\scriptsize $a x + b y + c z$};
  \foreach \l in {l1,l2,l3,l4}{
    \draw[refedge] (trunk.south) -- (\l.north);}
  \node[font=\scriptsize\itshape, text=simGray, align=center] at (0,-3.1)
    {fused accelerated kernels --- ``stopped low,'' specialization \emph{notes}\\
     the \emph{combinator} rises; the leaves are labels for fixed term-list lengths};
\end{tikzpicture}
\caption[\texttt{linear\_combination} and its arity leaves]{The combinator-primary
picture of the tensor-sum vocabulary. \code{linear\_combination} (trunk) is the
one combinator: a sum over a term list of any length. The four named BLAS routines
(leaves) are its fixed-arity specializations---\code{scal} at one term,
\code{axpy}/\code{axpby} at two, \code{axpbypcz} at three. They are
performance-fused kernels with no abstraction value of their own, so they are
``stopped low'' as specialization notes while the combinator rises to the top of
the vocabulary. One combinator, four leaves. \Cref{wex:arity} traces the
derivation with numbers.}
\label{fig:lincomb}
\end{figure}

\subsection{\texttt{inner\_product}: the scalar-producing fold, and its leaves}

The second data combinator is \code{inner\_product}, the scalar-producing sibling
of \code{linear\_combination}. Where \code{linear\_combination} folds a term list
into a \emph{tensor}, \code{inner\_product} folds two tensors (through an operator
$M$) into a \emph{scalar}, $\langle \vx, \vy\rangle_M = \vx^{\!\top} M \vy$. It is a
\emph{fold} in the same sense the running sum of \cref{ch:fold} was: it walks the
two operands position by position, accumulating a running scalar
$\sum_k (M\vy)_k\,\overline{x_k}$, and reduces the whole to a single number. The two
data combinators are thus the two halves of the L4 algebra of folds---one folds to
a tensor, the other folds to a scalar---and the step bodies of the iteration
combinators consume both: a Krylov step forms an \code{inner\_product} to compute a
step length and a \code{linear\_combination} to apply it.

\code{inner\_product} too carries a cluster of named leaves, but the cluster has a
subtlety worth drawing (\cref{fig:innerprod}).

\code{dot} is a genuine \emph{specialization}: it is \code{inner\_product} at the
special case $M = I$, the plain Euclidean (or Hermitian) inner product
$\vx^{\!\top}\vy$ that a solver spells \lstinline[style=l4style]|dot(p, Ap)|. Like
the BLAS arity leaves, it is the combinator at a fixed setting---here a fixed
operator argument rather than a fixed list length---and it hangs beneath the
combinator as a note.

\code{nrm2}, the Euclidean norm $\norm{\vx}_2 = \sqrt{\vx^{\!\top}\vx}$, is
\emph{not} a specialization of \code{inner\_product}. It is a \emph{consumer}: it
is built \emph{on top of} \code{dot} (itself a specialization), composing a square
root and an absolute value after the inner product, $\code{nrm2} =
\sqrt{\,\cdot\,} \circ \code{dot}$ at the diagonal. The distinction matters for the
combinator-primary bookkeeping: a \emph{leaf} is the combinator at a fixed setting
and lives beneath it; a \emph{consumer} calls the combinator and lives
\emph{above} it. Drawing \code{nrm2} as a leaf of \code{inner\_product} would
misfile a caller as a special case.

\begin{figure}[t]
\centering
\begin{tikzpicture}[font=\footnotesize, node distance=6mm]
  \node[stage, minimum width=50mm, align=center] (ip) at (0,0)
    {\code{inner\_product} \;\itshape\quad $\langle x, y\rangle_M = x^{\!\top} M y$};
  % dot: a specialization BELOW
  \node[opbox, minimum width=30mm, align=center] (dot) at (-1.6,-2.0)
    {\code{dot} \;\scriptsize (leaf)\\\scriptsize $M = I$: \;$x^{\!\top} y$};
  \draw[refedge] (ip.south) -- node[left,font=\scriptsize,pos=0.55]{\emph{specializes}} (dot.north);
  % nrm2: a CONSUMER above dot
  \node[product, minimum width=30mm, align=center] (nrm) at (-1.6,-4.0)
    {\code{nrm2} \;\scriptsize (consumer)\\\scriptsize $\sqrt{\;}\circ$ \code{dot}: \;$\|x\|_2$};
  \draw[depedge] (nrm.north) -- node[right,font=\scriptsize,pos=0.5]{\emph{consumes}} (dot.south);
  % annotation
  \node[font=\scriptsize\itshape, text=simGray, align=left, anchor=west] at (2.6,-2.0)
    {a \emph{leaf} is the combinator\\at a fixed setting ($M=I$):\\it lives \emph{below}};
  \node[font=\scriptsize\itshape, text=simGray, align=left, anchor=west] at (2.6,-4.0)
    {a \emph{consumer} \emph{calls} the\\combinator (here via \code{dot}):\\it lives \emph{above}};
\end{tikzpicture}
\caption[\texttt{inner\_product}: specialization versus consumer]{The
inner-product cluster, distinguishing two relationships the combinator-primary
model must keep separate. \code{dot} is a \emph{specialization} of
\code{inner\_product}---the combinator at the fixed operator $M = I$---and so hangs
\emph{below} it, exactly as the BLAS arity leaves hang below
\code{linear\_combination}. \code{nrm2} is not a special case but a \emph{consumer}:
it is built \emph{on} \code{dot}, composing a square root after the inner product,
and so lives \emph{above} it. A leaf is the combinator narrowed; a consumer is the
combinator called. Misdrawing \code{nrm2} as a leaf would file a caller as a
special case.}
\label{fig:innerprod}
\end{figure}

\begin{keyidea}
The data-algebra has two combinators. \code{linear\_combination} folds a term list
into a tensor, $\sum_i a_i t_i$; the BLAS routines \code{scal}, \code{axpy},
\code{axpby}, \code{axpbypcz} are its fixed-\emph{arity} leaves (fused kernels
stopped low, the combinator rising). \code{inner\_product} folds two tensors into a
scalar, $\langle x,y\rangle_M$; \code{dot} is its $M = I$ leaf, and \code{nrm2} is
a \emph{consumer} of \code{dot}, not a leaf. The pattern is uniform: \emph{one
combinator, several named specializations}, with a clean line between a
specialization (below) and a consumer (above).
\end{keyidea}

\begin{pitfall}
It is tempting to unify the two data combinators into one ``reduce a list of
tensors'' verb, since both are folds. Resist, for the same reason \cref{ch:situating}
kept the Gram and the projection reductions apart: sharing a \emph{shape} (both are
folds) is not sharing a \emph{meaning}. \code{linear\_combination} folds a list of
scalar-weighted tensors into a \emph{tensor} of the same shape; \code{inner\_product}
folds \emph{two} tensors through an operator into a \emph{scalar}. Their outputs
differ in rank, their inputs differ in arity, and their laws differ in kind
(\code{linear\_combination} is a monoid homomorphism over concatenation;
\code{inner\_product} is bilinear in its two arguments). Forcing them together would
either burden the tensor sum with a needless operator argument or strip the inner
product of its bilinearity. Two combinators that happen to both be folds is the
right design---the same judgment, applied to data verbs, that kept two reduction
verbs apart in \cref{ch:situating}.
\end{pitfall}

\section{The combinator-primary model}
\label{sec:primary}

We can now state the thesis the chapter has been assembling. Look back over the two
preceding sections and one shape recurs in every case. There is a \emph{combinator}
at the top---\code{solve\_family}, \code{fold\_solve}, \code{linear\_combination},
\code{inner\_product}---and beneath it a cluster of \emph{named operators} that are
each the combinator at some fixed setting: a fixed operator (\code{dot} at
$M = I$), a fixed term-list length (the BLAS arity leaves), a fixed family shape
(electrostatics as \code{solve\_family}). The combinator is the \emph{entry}; the
named operators are \emph{specialization notes} hung beneath it. This is the
\textbf{combinator-primary model}, and it is how the entire vocabulary of the
simulator is organized.

\begin{definition}[Combinator-primary organization]
\label{def:primary}
A vocabulary is \emph{combinator-primary} when its load-bearing entries are a small
set of higher-order \emph{combinators}---each carrying a signature, a set of
algebraic laws, and a structural meaning---and every other named operator is
recorded as a \emph{specialization} of some combinator: the combinator at a fixed
argument, a fixed arity, or a fixed family shape, written as a thin note that points
\emph{up} to its combinator and restates none of the combinator's structure. A
\emph{consumer} of a combinator (an operator that calls it) is distinguished from a
specialization (an operator that is the combinator narrowed) and is recorded
\emph{above} the combinator, not beneath it.
\end{definition}

\Cref{fig:primary} draws the model against its opposite. On the left, the
combinator-primary organization: one entry, many notes, every named operator
reachable as a setting of the combinator above it. On the right, the organization
the framework deliberately rejects: a flat pile of special-case kernels---
\code{scal}, \code{axpy}, \code{axpby}, \code{axpbypcz}, \code{dot}, \code{nrm2},
$n$ separate solver loops---each documented in full, with no unifying combinator
and so with the shared structure repeated $n$ times. The pile is not wrong, exactly;
it is \emph{unmanageable}. Every law must be restated per kernel; every reader must
rediscover that \code{axpby} and \code{axpbypcz} are the same idea at different
arities; every new special case adds a new full entry.

\begin{figure}[t]
\centering
\begin{tikzpicture}[font=\footnotesize]
  % ===== LEFT: combinator-primary =====
  \begin{scope}
    \node[font=\small\bfseries, text=simGreen] at (1.4,3.0) {combinator-primary (this book)};
    \node[stage, minimum width=40mm, align=center] (comb) at (1.4,2.0)
      {\textbf{combinator} \itshape (the entry)};
    \foreach \i/\x in {1/-0.5,2/1.4,3/3.3}{
      \node[opbox, minimum width=14mm, minimum height=6mm] (nt\i) at (\x,0.5) {};
      \draw[refedge] (comb.south) -- (nt\i.north);}
    \node[font=\scriptsize\itshape, text=simGray, align=center] at (1.4,-0.3)
      {specialization \emph{notes}\\(thin: a setting of the entry)};
  \end{scope}
  \draw[simGray!50, thick] (4.6,-0.7) -- (4.6,3.2);
  % ===== RIGHT: the pile =====
  \begin{scope}[xshift=5.6cm]
    \node[font=\small\bfseries, text=simRust] at (1.6,3.0) {the pile (rejected)};
    \foreach \i/\x/\y in {1/0/2.0,2/1.6/2.0,3/3.2/2.0,4/0/1.0,5/1.6/1.0,6/3.2/1.0}{
      \node[product, minimum width=13mm, minimum height=6mm] (pk\i) at (\x,\y) {};}
    \node[font=\scriptsize\itshape, text=simRust, align=center] at (1.6,-0.1)
      {$n$ full special-case kernels,\\no unifying combinator:\\shared structure repeated $n$ times};
  \end{scope}
\end{tikzpicture}
\caption[The combinator-primary model versus the pile]{The combinator-primary
model (left) and its rejected opposite (right). Combinator-primary: one combinator
is the \emph{entry}, carrying the signature, the laws, and the structure; the named
operators beneath it are thin specialization \emph{notes}, each a setting of the
combinator. The pile (right): a flat collection of special-case kernels, each a
full entry, with no combinator above to factor the shared structure---so the
structure is restated once per kernel and the collection grows without bound. The
left organization is why this book is short.}
\label{fig:primary}
\end{figure}

\subsection{Replace-and-propagate, not mine-and-strand}

The combinator-primary model is not just a way to \emph{present} the vocabulary;
it is a discipline for \emph{building} it. When a recurring pattern is spotted
across several operators---say, that \code{axpy} and \code{axpby} are both
two-term sums---there are two things one might do with the observation. The wrong
move is to \emph{mine and strand}: name the combinator, write it down in a corner,
and leave the original special-case kernels standing untouched beside it. Now there
are \emph{more} entries than before, not fewer: the kernels \emph{and} an unused
combinator. The right move is to \emph{replace and propagate}: make the combinator
the entry, rewrite each special case as a thin note that points up to it, and
propagate the change so that every consumer now reads through the combinator. The
count goes \emph{down}; the structure is stated \emph{once}.

\begin{keyidea}
\textbf{The combinator-primary model.} A small set of higher-order combinators are
the load-bearing vocabulary; each named operator is a thin \emph{specialization
note} beneath the combinator it instantiates. The combinator is the entry---it
carries the signature, the laws, the structure---and the leaves are settings of it.
The discipline is \emph{replace-and-propagate, not mine-and-strand}: a newly found
combinator \emph{replaces} the special cases (which become notes pointing up to it)
rather than standing as one more entry beside them. This is why the book can be
short: a few combinators and many specializations, with the shared structure stated
once rather than once per operator.
\end{keyidea}

\section{Why this matters downstream}
\label{sec:downstream}

The combinator-primary model is not an aesthetic preference; it is the load-bearing
beam of the rest of the book. Two later Parts rest on it directly, and it is worth
naming the dependence now so that those Parts read as payoffs rather than
surprises.

First, the \emph{drivers}. \Cref{part:modalities} studies the simulator's analyses
one at a time---electrostatic, magnetostatic, driven, transient, eigenmode,
boundary-mode. The combinator-primary model collapses what would be six
independent studies into six \emph{settings}. Every driver is one of
$\{\code{map}, \code{fold}, \code{single-call}\}$ over the
\texttt{assemble} $\to$ \texttt{solve} $\to$ \texttt{reduce} skeleton of
\cref{ch:situating}: electrostatics and magnetostatics are \code{solve\_family}
maps; the driven analysis is a \code{frequency\_sweep} map; the transient analysis
and the AMR loop are \code{fold\_solve} folds; the eigenmode and boundary-mode
analyses are single \code{eigsolve} calls. \Cref{fig:driverskel} previews this:
a driver is a combinator wrapped around the skeleton, with only the box contents
and the choice of combinator changing from analysis to analysis.

\begin{figure}[t]
\centering
\begin{tikzpicture}[font=\footnotesize, node distance=8mm]
  % the skeleton body
  \node[stage] (asm) at (0,0) {assemble};
  \node[stage, right=of asm] (solve) {solve};
  \node[stage, right=of solve] (red) {reduce};
  \draw[flow] (asm) -- node[above,font=\scriptsize]{$\Kop,\dots$} (solve);
  \draw[flow] (solve) -- node[above,font=\scriptsize]{field(s)} (red);
  % combinator wrapper
  \node[draw=simGreen, thick, rounded corners=3pt, fit=(asm)(solve)(red),
    inner sep=7mm, fill=simBgGreen, fill opacity=0.25] (wrap) {};
  \node[anchor=south west, font=\scriptsize\bfseries, text=simGreen]
    at (wrap.north west) {a combinator: \code{map} / \code{fold} / single-call};
  % the family/carry edge
  \draw[backflow, simGreen]
    (red.south) ++(0,-2mm) .. controls +(0,-7mm) and +(0,-7mm) ..
    (solve.south)
    node[midway, below, font=\scriptsize, text=simGreen]
      {next family member (map) \;or\; next carry (fold)};
  % in/out
  \node[data, left=10mm of wrap] (cfg) {config};
  \node[product, right=10mm of wrap] (prod) {product};
  \draw[flow] (cfg) -- (wrap.west);
  \draw[flow] (wrap.east) -- (prod);
\end{tikzpicture}
\caption[A driver as a combinator over the skeleton]{Every driver, previewed as a
combinator wrapped around the \texttt{assemble} $\to$ \texttt{solve} $\to$
\texttt{reduce} skeleton of \cref{ch:situating}. The combinator (green wrapper)
supplies the family structure---a \code{map} iterating over independent family
members, a \code{fold} threading a carry, or a single opaque call---and the
skeleton body supplies the per-member work. \Cref{part:modalities} fills in the box
contents one analysis at a time; what changes from driver to driver is only the box
contents and the choice of combinator. The combinator-primary model is what makes
six analyses six settings of this one picture.}
\label{fig:driverskel}
\end{figure}

To see how thoroughly the combinators carve up the six analyses, read down this
mapping---it is the punchline of \cref{ch:situating}'s grid restated in combinator
terms, and it is the skeleton of \cref{part:modalities}:
\begin{center}
\setlength{\tabcolsep}{8pt}
\renewcommand{\arraystretch}{1.2}
\begin{tabular}{@{}lll@{}}
\toprule
Driver & Combinator & Family structure \\
\midrule
Electrostatic & \code{solve\_family} & map over conductors (fixed $\Kop$) \\
Magnetostatic & \code{solve\_family} & map over current loops (fixed $\Kop$) \\
Driven & \code{frequency\_sweep} & map over $\omega$ (rebuild $\mat{A}(\omega)$) \\
Transient & \code{fold\_solve} & fold over timesteps (thread field state) \\
Eigenmode & \code{eigsolve} & single opaque call \\
Boundary mode & \code{eigsolve} & single opaque call (on a cross-section) \\
\bottomrule
\end{tabular}
\end{center}
Six drivers, four combinators, three structural shapes
($\code{map}$, $\code{fold}$, single-call). The physics that distinguishes the rows
lives in the \emph{box contents}---which operator is assembled, on which function
space, reduced by which verb---but the \emph{control structure} of each driver is
one of these four combinators wrapped around the skeleton. That is the entire reason
\cref{part:modalities} can present six analyses as six short variations rather than
six full programs.

Second, the \emph{synthesized library}. \Cref{part:collection} renders the whole
specification as though it were a real implementation library---the code the spec
describes, written out. That library \emph{is} these combinators. Its modules are
organized around them: an iteration module holding \code{iterate\_while} and the
step kernels, a data-algebra module holding \code{linear\_combination} and
\code{inner\_product} with their leaves, a coordination module holding
\code{solve\_family}, \code{frequency\_sweep}, and \code{fold\_solve}, and a driver
module composing them into the five analyses. The combinator-primary model is the
table of contents of that library. Learn the combinators here and you have learned
the spine of the synthesized code.

\section{Worked examples}

\begin{workedexample}{\texttt{solve\_family} as a map, \texttt{fold\_solve} as a foldl}{mapvsfold}
We write the two family combinators side by side on a concrete three-element
family and read off why one parallelizes and the other does not.

\emph{The map.} \code{solve\_family} runs three independent solves against one
fixed operator $\Kop$ and three right-hand sides $\vb_1, \vb_2, \vb_3$:
\begin{lstlisting}[style=l4style]
solve_family K [b1, b2, b3]
  = map (\b -> ksp_solve K b) [b1, b2, b3]
  = [ ksp_solve K b1            -- depends only on (K, b1)
    , ksp_solve K b2            -- depends only on (K, b2)
    , ksp_solve K b3 ]          -- depends only on (K, b3)
\end{lstlisting}
Each element names its own inputs $(\Kop, \vb_i)$ and \emph{nothing else}. No
element reads another element's output. So the three solves can be evaluated in any
order---$\vb_3$ first, $\vb_1$ last---or all at once on three machines, and the
collected list is the same either way. Concretely, with
$\Kop = \begin{psmallmatrix}2&0\\0&3\end{psmallmatrix}$ and
$\vb_1=\begin{psmallmatrix}2\\0\end{psmallmatrix},
\vb_2=\begin{psmallmatrix}0\\3\end{psmallmatrix},
\vb_3=\begin{psmallmatrix}2\\3\end{psmallmatrix}$, the three solves of
$\Kop\vx_i=\vb_i$ give $\vx_1=\begin{psmallmatrix}1\\0\end{psmallmatrix},
\vx_2=\begin{psmallmatrix}0\\1\end{psmallmatrix},
\vx_3=\begin{psmallmatrix}1\\1\end{psmallmatrix}$ independently---and indeed
$\vx_3 = \vx_1 + \vx_2$ only because $\vb_3 = \vb_1 + \vb_2$ and the solves are
linear, not because element three \emph{read} elements one and two. The map never
threads a carry.

\emph{The fold.} \code{fold\_solve} runs three \emph{chained} steps, threading a
state $s$ through a schedule $[t_1, t_2, t_3]$:
\begin{lstlisting}[style=l4style]
fold_solve op s0 [t1, t2, t3]
  = foldl (\s t -> time_step_op op s t) s0 [t1, t2, t3]
  = time_step_op op
      (time_step_op op
         (time_step_op op s0 t1)   -- produces s1
         t2)                       -- s2 needs s1
      t3                           -- s3 needs s2
\end{lstlisting}
Now read the nesting from the inside out. The innermost call produces $s_1$ from
$s_0$; the middle call needs $s_1$ to produce $s_2$; the outer call needs $s_2$ to
produce $s_3$. The step at $t_2$ \emph{cannot start} until the step at $t_1$ has
finished, because $s_1$ is literally its input. There is no order but
$t_1, t_2, t_3$, and no machine can take the second step early. With a toy
$\code{time\_step\_op}\ op\ s\ t = s + 1$ and $s_0 = 0$ the fold computes
$0 \to 1 \to 2 \to 3$---each value one more than the last, and unreachable except
by passing through its predecessor.

\emph{The contrast.} The two pieces of pseudocode differ by exactly one structural
fact: the map's element function names only its own inputs, while the fold's step
function names the prior carry $s$. That single carry dependence is the difference
between a fan-out and a chain (\cref{fig:mapfold}), between parallel and
sequential. It is why electrostatics (a \code{solve\_family} map) can solve all its
conductors at once while the transient analysis (a \code{fold\_solve} fold) must
march its timesteps in order. The map is the degenerate fold whose step would
ignore $s$; the fold is the one that uses it.
\end{workedexample}

\begin{workedexample}{One combinator, four leaves: deriving the BLAS family}{arity}
We take \code{linear\_combination} and recover each named BLAS routine as a fixed-
arity specialization, with numbers, so the ``one combinator, four leaves'' claim is
concrete rather than slogan. Throughout, let
$\vx = \begin{psmallmatrix}1\\2\end{psmallmatrix}$,
$\vy = \begin{psmallmatrix}3\\4\end{psmallmatrix}$,
$\vz = \begin{psmallmatrix}5\\6\end{psmallmatrix}$, and coefficients
$a = 2,\ b = 10,\ c = 100$.

\emph{Arity 1 ---} \code{scal}. The one-term list:
\begin{lstlisting}[style=l4style]
scal(a, x) = linear_combination [(a, x)] = zeros + scal a x = a*x
\end{lstlisting}
Numerically $\code{scal}(2, \vx) = 2\vx = \begin{psmallmatrix}2\\4\end{psmallmatrix}$.
The fold seeds at \code{zeros} and folds one term: $0 + 2\vx$.

\emph{Arity 2, second coefficient fixed ---} \code{axpy}. The two-term list whose
second coefficient is $1$:
\begin{lstlisting}[style=l4style]
axpy(a, x, y) = linear_combination [(a, x), (1, y)] = a*x + 1*y = a*x + y
\end{lstlisting}
Numerically $\code{axpy}(2, \vx, \vy) = 2\vx + \vy =
\begin{psmallmatrix}2\\4\end{psmallmatrix} + \begin{psmallmatrix}3\\4\end{psmallmatrix}
= \begin{psmallmatrix}5\\8\end{psmallmatrix}$.

\emph{Arity 2, general ---} \code{axpby}. The two-term list with both coefficients
free:
\begin{lstlisting}[style=l4style]
axpby(a, x, b, y) = linear_combination [(a, x), (b, y)] = a*x + b*y
\end{lstlisting}
Numerically $\code{axpby}(2, \vx, 10, \vy) = 2\vx + 10\vy =
\begin{psmallmatrix}2\\4\end{psmallmatrix} + \begin{psmallmatrix}30\\40\end{psmallmatrix}
= \begin{psmallmatrix}32\\44\end{psmallmatrix}$. Note \code{axpy} is the special
case $b = 1$: $\code{axpby}(a, \vx, 1, \vy) = \code{axpy}(a, \vx, \vy)$.

\emph{Arity 3 ---} \code{axpbypcz}. The three-term list:
\begin{lstlisting}[style=l4style]
axpbypcz(a, x, b, y, c, z) = linear_combination [(a, x), (b, y), (c, z)]
                           = a*x + b*y + c*z
\end{lstlisting}
Numerically $\code{axpbypcz}(2,\vx,10,\vy,100,\vz) = 2\vx + 10\vy + 100\vz =
\begin{psmallmatrix}2\\4\end{psmallmatrix}+\begin{psmallmatrix}30\\40\end{psmallmatrix}
+\begin{psmallmatrix}500\\600\end{psmallmatrix}=\begin{psmallmatrix}532\\644\end{psmallmatrix}$.

\emph{The concatenation law makes the subsumption exact.} Watch the three-term sum
split into a two-term and a one-term list:
\[
  \underbrace{\code{linear\_combination}\,[(a,\vx),(b,\vy),(c,\vz)]}_{\code{axpbypcz}}
  \;=\;
  \underbrace{\code{linear\_combination}\,[(a,\vx),(b,\vy)]}_{\code{axpby}}
  \;+\;
  \underbrace{\code{linear\_combination}\,[(c,\vz)]}_{\code{scal}}.
\]
Numerically $\begin{psmallmatrix}532\\644\end{psmallmatrix} =
\begin{psmallmatrix}32\\44\end{psmallmatrix} + \begin{psmallmatrix}500\\600\end{psmallmatrix}$:
the \code{axpbypcz} result is the \code{axpby} result plus the \code{scal} result.
The four named routines are not four ideas; they are one combinator read off at
list-lengths one, two, and three, and the concatenation-homomorphism is the law
that stitches them together. This is the combinator-primary model in a single
trace: the combinator carries the structure (the law), and the leaves are settings
of it.
\end{workedexample}

\summarysection
\Cref{ch:fold} built \code{iterate\_while}, the one iteration primitive; this
chapter widened the lens to the combinators that organize the whole simulator. The
governing axis is \emph{map versus fold}: a map runs a family of \emph{independent}
computations (no carry, so they parallelize), while a fold threads a \emph{carry}
through a chain (each step needs the last, so it serializes); a map is the
degenerate fold whose step ignores the carry. The four solve corners of
\cref{ch:situating} are exactly four combinators: \code{solve\_family} (the
fixed-operator map; electrostatic, magnetostatic), \code{frequency\_sweep} (the
operator-varying map, rebuilding the operator inside the loop; driven),
\code{fold\_solve} (the state-threaded fold; transient and AMR), and \code{eigsolve}
(the opaque black-box, no loop we write; eigenmode, boundary-mode). The
data-algebra adds two combinators: \code{linear\_combination}, the tensor-sum
$\sum_i a_i t_i$, whose fixed-arity leaves are the BLAS routines
\code{scal}/\code{axpy}/\code{axpby}/\code{axpbypcz} (fused kernels stopped low,
the combinator rising); and \code{inner\_product}, the scalar-producing fold, with
\code{dot} as its $M = I$ leaf and \code{nrm2} as a \emph{consumer} of \code{dot},
not a leaf. The thesis is the \emph{combinator-primary model}: the combinator is
the entry, carrying the signature and laws; the named operators are thin
specialization notes beneath it; the discipline is \emph{replace-and-propagate, not
mine-and-strand}. This is why a book about a vast simulator can be short---a few
combinators, many specializations---and it is the spine of both the drivers
(\cref{part:modalities}, each a combinator over assemble--solve--reduce) and the
synthesized library (\cref{part:collection}, which \emph{is} these combinators).

\furtherreading
The \code{map}/\code{fold} vocabulary, and the laws that make a fold-of-fixed-length
the same as its general fold, are foundational functional programming; Peyton
Jones's account of Haskell~\citep{peytonjones2003} develops \code{map}, \code{foldl},
and the homomorphism laws this chapter leans on for the arity leaves. The BLAS
routines whose fused special cases sit beneath \code{linear\_combination} and
\code{inner\_product}---the \code{axpy} family, \code{dot}, \code{nrm2}---and the
reason they are written as fused kernels rather than general sums are documented in
the numerical-linear-algebra literature; Golub and Van Loan~\citep{golubvanloan2013}
is the standard reference for the BLAS-1 level and its performance model, and
Saad~\citep{saad2003} shows how Krylov solvers are spelled in exactly these verbs,
which is the form \cref{ch:cg} and \cref{ch:gmres} build on.

\exercisessection
\begin{enumerate}[nosep]
  \item For each of the four solve corners, name its combinator and state the two
  binary classifications that place it (operator fixed/varying; map/fold or
  neither). Which corner answers ``neither row'' of the map/fold axis, and why?
  \item Electrostatics is \code{solve\_family} and the driven analysis is
  \code{frequency\_sweep}. Both are maps. Pinpoint the single line of pseudocode
  whose position---inside the map versus hoisted outside it---is the entire
  difference between them, and explain why moving it changes the \emph{answer}, not
  merely the speed.
  \item Write \code{axpby} and \code{axpy} as instances of \code{linear\_combination},
  then show that \code{axpy} is the special case of \code{axpby} with one
  coefficient fixed. Using the concatenation-homomorphism, express \code{axpbypcz}
  as a sum of an \code{axpby} and a \code{scal}, and verify it numerically with
  $\vx=\begin{psmallmatrix}1\\0\end{psmallmatrix},
  \vy=\begin{psmallmatrix}0\\1\end{psmallmatrix},
  \vz=\begin{psmallmatrix}1\\1\end{psmallmatrix}$ and coefficients $1, 2, 3$.
  \item Explain why \code{dot} is a \emph{specialization} of \code{inner\_product}
  but \code{nrm2} is a \emph{consumer} of \code{dot}, not a specialization of
  \code{inner\_product}. In the combinator-primary picture, which one hangs
  \emph{below} its combinator and which sits \emph{above}? What would go wrong if
  you drew \code{nrm2} as a leaf of \code{inner\_product}?
  \item Contrast \emph{replace-and-propagate} with \emph{mine-and-strand}. A
  colleague notices that \code{axpy} and \code{axpby} are both two-term sums, writes
  down a general \code{two\_term} combinator in a corner, and leaves the original
  \code{axpy} and \code{axpby} entries untouched. Has the entry count gone up or
  down? Which discipline did they follow, and what should they have done instead?
  \item A \code{map} is ``the degenerate fold whose step ignores the carry.'' Write
  \code{map f xs} as a \code{fold\_solve}-style \code{foldl} whose step discards its
  accumulator. Then argue from your rewrite why a map parallelizes and a fold does
  not---i.e., which data dependence the discarded carry would have introduced.
  \item Using \cref{fig:driverskel}, describe the transient analysis and the
  electrostatic analysis as ``a combinator wrapped around assemble--solve--reduce.''
  Which combinator does each use, what flows along the back-edge in each case (a
  family member or a threaded carry?), and which of the two could you run on a
  cluster of independent machines?
  \item \emph{(Synthesis.)} The chapter claims the book is short \emph{because} of
  the combinator-primary model. Suppose instead the simulator's vocabulary were
  written as ``the pile'' of \cref{fig:primary}: six independent driver studies,
  six special-case family loops, four BLAS kernels, two inner-product kernels, each
  a full entry. Estimate, in rough terms, how the page count would grow, and name
  two specific facts (one law, one structural relationship) that would have to be
  restated multiple times rather than once.
\end{enumerate}
